\section{Responsibility-targeted reframing and dependency-directed reopening}

\subsection{Why reframing is dangerous}

The ability to rewrite a problem is useful precisely because it is high impact. An unconstrained system can explain every failure by changing the question, the representation, or the method until the failure disappears. Paper I therefore treats reframing as a licensed state transition rather than a generic recovery heuristic.

Let $r_t$ be a typed residual and $d_t$ the current diagnosis. A reframe is admissible only when the evidence carried by $d_t$ supports a formulation/search responsibility coordinate. Conceptually,
\[
\mathrm{REFRAME}(X_t,r_t,d_t)
\rightarrow X_{t+1}
\]
is partial: unsupported diagnoses must return no formulation change. The local implementation now rejects the shortcut
\[
\text{singular responsibility}\Rightarrow\text{reframe}.
\]
Singularity is not enough. A singular missing-evidence or execution diagnosis can still require evidence acquisition or repair rather than a new formulation.

\subsection{Scoped reopening}

Suppose a conclusion $c$ was closed under dependencies $D(c)$. When a reframe changes coordinate $z$, ORION reopens $c$ only if $z$ is in the transitive dependency basis of the closure:
\[
z\in D^{*}(c)\Rightarrow \mathrm{status}(c)\leftarrow \mathrm{STALE}.
\]
This rule avoids two opposite errors. Failing to reopen produces stale certainty: old conclusions remain ``solved'' after their assumptions changed. Reopening the entire project after every change destroys locality and makes recursive research unnecessarily expensive.

The key evaluation object is therefore not only final task success. Reopening has its own precision/recall question: did the system stale exactly the closure that materially depended on the changed coordinate?

\subsection{Bounded recursion, not completeness}

Paper I does not claim that recursion reaches an open-world complete state. The stopping condition is operational and scoped: continue reconstruction while material residuals or affected reopen obligations remain, subject to the stopping and coverage rules delegated to Paper II. A bounded stop means only that the registered attacks and obligations are flat at the declared cutoff. A new representation, domain, residual, or formulation change can reopen the process.

This distinction prevents ``nothing new was found'' from being promoted into ``nothing relevant exists.'' It also keeps Paper I focused: the scientific question here is whether explicit reconstruction and scoped reopening help under hidden formulation shifts, not whether the search procedure achieves exhaustive recall.
